
\section{Installation}
VSPLINE
is distributed as a Bourne Shell ``bundle''.
Create a directory called {\tt vspline},
move the bundled file to that directory,
and type {\tt sh} {\em filename} at a UNIX prompt.
The file {\tt README} contains important
installation instructions.

When installed,
the C source code resides in several subdirectories.
{\tt matrix} contains very general vector
and matrix utility subroutines.
Simply by changing the definitions in
{\tt matrix/matrix.h},
one could use a different library.
{\tt smooth}
contains utilities for banded matrices,
and the smoothing subroutines themselves.
Finally,
{\tt Test}
contains programs we used to test the library.
The test programs were applied to
the data files in {\tt Data};
to verify your installation,
follow the instructions in {\tt Data/README}.
One {\tt Makefile} in the parent directory
controls compilation of all three subdirectories.
This Makefile creates five libraries:
{\tt libnonlin.a, libsmooth.a,
liboptima.a, libmatrix.a, libiodep.a},
which should be linked {\em in the order given}
when compiling.
The Makefile in the Test directory provides a good
example.

%The {\tt powell} subroutine
%of the Numerical Recipes in C library~\cite{numrecc:88}
%is used to minimize the CV score.
%Therefore, you must load in both {\tt libvspline.a}
%and (e.g.) {\tt libnumrecC.a}.
%If you do not have this library,
%type {\tt make libvsmooth.a}.
%This latter library is a subset of
%the VSPLINE library
%that contains the basic smoothing routines
%(sufficient to use the {\tt vspline} routine),
%but does not contain the routines
%for minimizing the CV score.

{\tt vspline} resides in {\tt smooth/example1.c},
{\tt bestcv} in {\tt smooth/example2.c}.
A subroutine called {\tt vect\_smooth}
in {\tt smooth/example1.c}
gives another example of how the VSPLINE library can used.
{\tt cubic\_interp} resides in {\tt smooth/interp.c}.
The \LaTeX source for this document is in {\tt smooth/Doc}.
